% !TeX program = xelatex
% vim: set textwidth=120:

% Example CV based on the 1.5-column-cv template. Main features:
% * uses the Roboto font family and IcoMoon icon set;
% * doesn't use colours, different font weights are used instead for styling;
% * because the CV fits on one page, header and footer is empty, since there isn't much useful info to put there;
% * includes a photo.
\documentclass[a4paper,10pt]{article}


% package imports
% ---------------

\usepackage[ngerman]{babel} % deutsche Trennung und Bezeichnungen
\usepackage{calc}           % for easier length calculations (infix notation)
\usepackage{enumitem}       % for configuring list environments
\usepackage{fancyhdr}       % for setting header and footer
\usepackage{fontspec}       % for fonts
\usepackage{geometry}       % for setting margins (\newgeometry)
\usepackage{graphicx}       % for pictures
\usepackage{microtype}      % for microtypography stuff
\usepackage{xcolor}         % for colours


% margin and column widths
% ------------------------

% margins
\newgeometry{left=15mm,right=15mm,top=15mm,bottom=15mm}

% width of the gap between left and right column
\newlength{\cvcolumngapwidth}
\setlength{\cvcolumngapwidth}{3.5mm}

% left column width
\newlength{\cvleftcolumnwidth}
\setlength{\cvleftcolumnwidth}{44.5mm}

% right column width
\newlength{\cvrightcolumnwidth}
\setlength{\cvrightcolumnwidth}{\textwidth-\cvleftcolumnwidth-\cvcolumngapwidth}

% set paragraph indentation to 0, because it screws up the whole layout otherwise
\setlength{\parindent}{0mm}


% style definitions
% -----------------
% style categories explanation:
% * \cvnameXXX is used for the name;
% * \cvsectionXXX is used for section names (left column, accompanied by a horizontal rule);
% * \cvtitleXXX is used for job/education titles (right column);
% * \cvdurationXXX is used for job/education durations (left column);
% * \cvheadingXXX is used for headings (left column);
% * \cvmainXXX (and \setmainfont) is used for main text;
% * \cvruleXXX is used for the horizontal rules denoting sections.

% font families
\defaultfontfeatures{Ligatures=TeX} % reportedly a good idea, see https://tex.stackexchange.com/a/37251

\newfontfamily{\cvnamefont}{Roboto Medium}
\newfontfamily{\cvsectionfont}{Roboto Medium}
\newfontfamily{\cvtitlefont}{Roboto Regular}
\newfontfamily{\cvdurationfont}{Roboto Light Italic}
\newfontfamily{\cvheadingfont}{Roboto Regular}
\setmainfont{Roboto Light}

% colours
\definecolor{cvnamecolor}{HTML}{000000}
\definecolor{cvsectioncolor}{HTML}{000000}
\definecolor{cvtitlecolor}{HTML}{000000}
\definecolor{cvdurationcolor}{HTML}{000000}
\definecolor{cvheadingcolor}{HTML}{000000}
\definecolor{cvmaincolor}{HTML}{000000}
\definecolor{cvrulecolor}{HTML}{000000}

\color{cvmaincolor}

% styles
\newcommand{\cvnamestyle}[1]{{\Large\cvnamefont\textcolor{cvnamecolor}{#1}}}
\newcommand{\cvsectionstyle}[1]{{\normalsize\cvsectionfont\textcolor{cvsectioncolor}{#1}}}
\newcommand{\cvtitlestyle}[1]{{\large\cvtitlefont\textcolor{cvtitlecolor}{#1}}}
\newcommand{\cvdurationstyle}[1]{{\small\cvdurationfont\textcolor{cvdurationcolor}{#1}}}
\newcommand{\cvheadingstyle}[1]{{\normalsize\cvheadingfont\textcolor{cvheadingcolor}{#1}}}


% inter-item spacing
% ------------------

% vertical space after personal info and standard CV items
\newlength{\cvafteritemskipamount}
\setlength{\cvafteritemskipamount}{5mm plus 1.25mm minus 1.25mm}

% vertical space after sections
\newlength{\cvaftersectionskipamount}
\setlength{\cvaftersectionskipamount}{2mm plus 0.5mm minus 0.5mm}

% extra vertical space to be used when a section starts with an item with a heading (e.g. in the skills section),
% so that the heading does not follow the section name too closely
\newlength{\cvbetweensectionandheadingextraskipamount}
\setlength{\cvbetweensectionandheadingextraskipamount}{1mm plus 0.25mm minus 0.25mm}


% intra-item spacing
% ------------------

% vertical space after name
\newlength{\cvafternameskipamount}
\setlength{\cvafternameskipamount}{3mm plus 0.75mm minus 0.75mm}

% vertical space after personal info lines
\newlength{\cvafterpersonalinfolineskipamount}
\setlength{\cvafterpersonalinfolineskipamount}{2mm plus 0.5mm minus 0.5mm}

% vertical space after titles
\newlength{\cvaftertitleskipamount}
\setlength{\cvaftertitleskipamount}{1mm plus 0.25mm minus 0.25mm}

% value to be used as parskip in right column of CV items and itemsep in lists (same for both, for consistency)
\newlength{\cvparskip}
\setlength{\cvparskip}{0.5mm plus 0.125mm minus 0.125mm}

% set global list configuration (use parskip as itemsep, and no separation otherwise)
\setlist{parsep=0mm,topsep=0mm,partopsep=0mm,itemsep=\cvparskip}


% CV commands
% -----------

% creates a "personal info" CV item with the given left and right column contents, with appropriate vertical space after
% @param #1 left column content (should be the CV photo)
% @param #2 right column content (should be the name and personal info)
\newcommand{\cvpersonalinfo}[2]{
    % left and right column
    \begin{minipage}[t]{\cvleftcolumnwidth}
        \vspace{0mm} % XXX hack to align to top, see https://tex.stackexchange.com/a/11632
        \raggedleft #1
    \end{minipage}% XXX necessary comment to avoid unwanted space
    \hspace{\cvcolumngapwidth}% XXX necessary comment to avoid unwanted space
    \begin{minipage}[t]{\cvrightcolumnwidth}
        \vspace{0mm} % XXX hack to align to top, see https://tex.stackexchange.com/a/11632
        #2
    \end{minipage}

    % space after
    \vspace{\cvafteritemskipamount}
}

% typesets a name, with appropriate vertical space after
% @param #1 name text
\newcommand{\cvname}[1]{
    % name
    \cvnamestyle{#1}

    % space after
    \vspace{\cvafternameskipamount}
}

% typesets a line of personal info beginning with an icon, with appropriate vertical space after
% @param #1 parameters for the \includegraphics command used to include the icon
% @param #2 icon filename
% @param #3 line text
\newcommand{\cvpersonalinfolinewithicon}[3]{
    % icon, vertically aligned with text (see https://tex.stackexchange.com/a/129463)
    \raisebox{.5\fontcharht\font`E-.5\height}{\includegraphics[#1]{#2}}
    % text
    #3

    % space after
    \vspace{\cvafterpersonalinfolineskipamount}
}

% creates a "section" CV item with the given left column content, a horizontal rule in the right column, and with
% appropriate vertical space after
% @param #1 left column content (should be the section name)
\newcommand{\cvsection}[1]{
    % left and right column
    \begin{minipage}[t]{\cvleftcolumnwidth}
        \raggedleft\cvsectionstyle{#1}
    \end{minipage}% XXX necessary comment to avoid unwanted space
    \hspace{\cvcolumngapwidth}% XXX necessary comment to avoid unwanted space
    \begin{minipage}[t]{\cvrightcolumnwidth}
        \textcolor{cvrulecolor}{\rule{\cvrightcolumnwidth}{0.3mm}}
    \end{minipage}

    % space after
    \vspace{\cvaftersectionskipamount}
}

% creates a standard, multi-purpose CV item with the given left and right column contents, parskip set to cvparskip
% in the right column, and with appropriate vertical space after
% @param #1 left column content
% @param #2 right column content
\newcommand{\cvitem}[2]{
    % left and right column
    \begin{minipage}[t]{\cvleftcolumnwidth}
        \raggedleft #1
    \end{minipage}% XXX necessary comment to avoid unwanted space
    \hspace{\cvcolumngapwidth}% XXX necessary comment to avoid unwanted space
    \begin{minipage}[t]{\cvrightcolumnwidth}
        \setlength{\parskip}{\cvparskip} #2
    \end{minipage}

    % space after
    \vspace{\cvafteritemskipamount}
}

% typesets a title, with appropriate vertical space after
% @param #1 title text
\newcommand{\cvtitle}[1]{
    % title
    \cvtitlestyle{#1}

    % space after
    \vspace{\cvaftertitleskipamount}
    % XXX need to subtract cvparskip here, because it is automatically inserted after the title "paragraph"
    \vspace{-\cvparskip}
}


% header and footer
% -----------------

% set empty header and footer
\pagestyle{empty}



% preamble end/document start
% ===========================

\begin{document}


% personal info
% -------------

\cvpersonalinfo{
    % Optional: Hier kann später ein Foto ergänzt werden.
}{
    % name
    \cvname{Christian Komposch}

    % address
    Echinger Str. 6b\\
    85386 Eching
    \vspace{\cvafterpersonalinfolineskipamount}

    % phone number
    +49-173-8482319
    \vspace{\cvafterpersonalinfolineskipamount}

    % email address
    christian@komposch.de
}


% work experience
% ---------------

\cvsection{BERUFSERFAHRUNG}

% Aktuelle Position
\cvitem{
    \cvdurationstyle{2026 -- heute}
}{
    \cvtitle{Produktverantwortlicher Software DCDC-Wandler und OBC}

    BMW AG

    \begin{itemize}[leftmargin=*]
        \item Fachliche Gesamtverantwortung für die Softwareroadmaps der DCDC-Wandler- und OBC-Steuergeräte
        \item Priorisierung und Steuerung des Produkt-Backlogs in enger Abstimmung mit Entwicklung, Qualität und Kunden
        \item % TODO: konkrete Kennzahl ergänzen, z. B. Anzahl betreuter Steuergerätevarianten/Fahrzeugprojekte, Teamgröße
              Verantwortung für [Anzahl] Steuergerätevarianten über [Anzahl] Fahrzeugprojekte hinweg
    \end{itemize}
}

\cvitem{
    \cvdurationstyle{2021 -- 2026}
}{
    \cvtitle{Produktverantwortlicher Software Batteriemanagementsystem}

    BMW AG

    \begin{itemize}[leftmargin=*]
        \item Fachliche Gesamtverantwortung für die Softwareroadmap des Batteriemanagementsystems
        \item Priorisierung und Steuerung des Produkt-Backlogs in enger Abstimmung mit Entwicklung, Qualität und Kunden
        \item 2025: Erfolgreicher Fahrzeuganlauf der Neuen Klasse (BMW iX3, NA5) als Inhouse-Entwicklungsprojekt
        \item % TODO: konkrete Kennzahl ergänzen, z. B. Teamgröße, Anzahl Lieferanten/Stakeholder
              Fachliche Führung eines Teams von [Anzahl] Entwicklern/Stakeholdern
    \end{itemize}
}

\cvitem{
    \cvdurationstyle{2016 -- 2021}
}{
    \cvtitle{Projektleiter Software Batteriemanagementsystem}

    BMW AG

    \begin{itemize}[leftmargin=*]
        \item Fachliche und organisatorische Leitung des Softwareentwicklungsteams für das Batteriemanagementsystem
        \item Verantwortung für Termin-, Budget- und Qualitätsziele über den gesamten Projektverlauf
        \item Steuerung und Koordination der beteiligten Lieferanten
        \item 2021: Erfolgreicher Fahrzeuganlauf des BMW iX (I20) -- termingerecht und in geforderter Qualität
        \item % TODO: konkrete Kennzahl ergänzen, z. B. Teamgröße, Budgetvolumen
              Verantwortung für ein Team von [Anzahl] Personen und ein Budget von [Betrag]
    \end{itemize}
}

\cvitem{
    \cvdurationstyle{01/2013 -- 2016}
}{
    \cvtitle{Softwareintegrator Embedded Systems / AUTOSAR -- Batteriemanagementsystem}

    BMW AG

    \begin{itemize}[leftmargin=*]
        \item Integration von Software für das Batteriemanagementsystem-Steuergerät (Embedded Systems / AUTOSAR)
        \item Koordination der Schnittstellen zwischen Entwicklungsteams zur Sicherstellung der Systemkonsistenz
        \item Tätigkeit im Bereich elektrischer Antriebe
        \item 2015: Fahrzeuganlauf des BMW X5 (G05)
        \item 2013: Fahrzeuganlauf des BMW i3 (I01)
    \end{itemize}
}

\cvitem{
    \cvdurationstyle{04/2012 -- 01/2013}
}{
    \cvtitle{Softwareintegrator Embedded Systems / AUTOSAR -- Hybridkomponenten}

    BMW Peugeot Citroën Electrification\\
    \textit{(Joint Venture der BMW AG mit PSA Peugeot Citroën)}

    \begin{itemize}[leftmargin=*]
        \item Integration von Software für ein Embedded-Systems-AUTOSAR-Steuergerät im Bereich Hybridantriebe
        \item Koordination der Schnittstellen zwischen Entwicklungsteams zur Sicherstellung der Systemkonsistenz
        \item Nahtlose Rückführung der Tätigkeit in die BMW AG nach Auflösung des Joint Ventures Anfang 2013
    \end{itemize}
}

\cvitem{
    \cvdurationstyle{09/2009 -- 04/2012}
}{
    \cvtitle{Berater Steuergeräteentwicklung / Embedded Systems}

    Altran Deutschland, im Auftrag der BMW AG

    \begin{itemize}[leftmargin=*]
        \item Entwicklung und Pflege einer Werkzeugkette zur automatischen Codegenerierung mit MATLAB/Simulink
        \item Umsetzung von Automatisierungslösungen zur Steigerung der Entwicklungseffizienz
        \item Entwicklung für ein Antriebssteuergerät
        \item Nach Bewährung in der Beraterrolle Übernahme in Festanstellung bei der BMW AG ab 2013
    \end{itemize}
}

\cvitem{
    \cvdurationstyle{01/2007 -- 09/2009}
}{
    \cvtitle{IT-Freelancer / IT-Berater}

    Selbständig unter eigenem Namen

    \begin{itemize}[leftmargin=*]
        \item Beratung und Betreuung kleiner und mittelständischer Unternehmen bei IT-Anliegen
        \item Technische Arbeiten in den Bereichen Server, Netzwerke, IT-Anwendungen, Datenbanken und Applikationsentwicklung
    \end{itemize}
}


% education
% ---------

\cvsection{AUSBILDUNG}

\cvitem{
    \cvdurationstyle{2004 -- 2009}
}{
    \cvtitle{Diplom-Ingenieur (FH) Softwaretechnik}

    Hochschule Esslingen

    \begin{itemize}[leftmargin=*]
        \item Abschlussnote: 2,3
        \item Diplomarbeit bei der IT-Designers GmbH, Esslingen, in Kooperation mit der Daimler AG, Sindelfingen: frühe
              Spezialisierung auf AUTOSAR -- Inbetriebnahme eines Vorserien-AUTOSAR-Stacks der Firma Vector Informatik
              sowie Erstellung einer prototypischen Anwendung nach dem AUTOSAR-Vorgehensmodell (Werkzeugkette DaVinci)
        \item Note der Diplomarbeit: 1,3
    \end{itemize}
}

\cvitem{
    \cvdurationstyle{2003 -- 2004}
}{
    \cvtitle{Fachhochschulreife}

    Friedrich-Ebert-Schule, Esslingen am Neckar

    \begin{itemize}[leftmargin=*]
        \item Abschlussnote: 2,2
    \end{itemize}
}

\cvitem{
    \cvdurationstyle{2000 -- 2003}
}{
    \cvtitle{Kommunikationselektroniker}

    Fachrichtung Informationstechnik

    Robert Bosch GmbH, Wernau

    \begin{itemize}[leftmargin=*]
        \item Abschlussnote: 1,8
    \end{itemize}
}


% skills
% ------

\cvsection{KENNTNISSE}

\vspace{\cvbetweensectionandheadingextraskipamount}

\cvitem{
    \cvheadingstyle{Sprachen}
}{
    Deutsch -- Muttersprache

    Englisch -- Verhandlungssicher
}

\cvitem{
    \cvheadingstyle{IT-Kenntnisse}
}{
    Entwicklungswerkzeuge -- GitHub, Zuul, Bazel, MATLAB/Simulink, MagicDraw

    Anforderungs-/Prozessmanagement -- Jira, Confluence, DOORS, Codebeamer

    Methoden \& Standards -- Softwareentwicklung, AUTOSAR, ASPICE, ISO 26262
}

\cvitem{
    \cvheadingstyle{Weitere Kenntnisse}
}{
    Fachkraft für Hochvolt-Systeme in Kraftfahrzeugen nach DGUV Information 209-093

    Führerschein Klasse B
}


% additional info
% ---------------

\cvsection{ZUSÄTZLICHE ANGABEN}

\vspace{\cvbetweensectionandheadingextraskipamount}

\cvitem{
    \cvheadingstyle{Interessen}
}{
    Radfahren, Sport, Heimvernetzung / Homelab-Projekte
}

\end{document}